% ---
% title: C++ static_assert 编译期断言
% date: 2026-07-13
% author: Crystal-Sky
% tags: C++, C++17, 基础语法, static_assert, 编译期检查
% summary: 对比 assert 与 static_assert，介绍编译期断言的语法、简化写法、类型检查和常见错误。
% slug: cpp-static-assert
% course: C++
% course_slug: cpp
% chapter: 基础语法
% chapter_order: 1
% lesson_order: 6
% lesson_type: 基础语法
% ---

\documentclass[UTF8,12pt,a4paper]{ctexart}
\usepackage{geometry}
\usepackage{xcolor}
\usepackage{listings}
\usepackage{booktabs}
\usepackage{enumitem}
\usepackage{tcolorbox}
\usepackage{fancyhdr}
\geometry{left=2.4cm,right=2.4cm,top=2.5cm,bottom=2.5cm}
\pagestyle{fancy}
\fancyhf{}
\fancyhead[L]{C++ static\_assert}
\fancyhead[R]{教学讲义}
\fancyfoot[C]{\thepage}
\definecolor{codebg}{RGB}{247,248,250}
\definecolor{keywordcolor}{RGB}{0,92,197}
\definecolor{commentcolor}{RGB}{0,128,0}
\definecolor{stringcolor}{RGB}{163,21,21}
\lstset{language=C++,basicstyle=\ttfamily\small,keywordstyle=\color{keywordcolor}\bfseries,commentstyle=\color{commentcolor},stringstyle=\color{stringcolor},backgroundcolor=\color{codebg},frame=single,numbers=left,numberstyle=\tiny\color{gray},numbersep=8pt,tabsize=4,showstringspaces=false,breaklines=true,columns=fullflexible,keepspaces=true}
\newtcolorbox{conceptbox}[1]{colback=blue!3,colframe=blue!60!black,title=#1,fonttitle=\bfseries}
\title{\textbf{C++ \texttt{static\_assert} 编译期断言}}
\author{}
\date{}
\begin{document}
\maketitle

\section{先回顾普通 \texttt{assert}}
普通 \texttt{assert} 用于在程序运行时检查条件。
\begin{lstlisting}
#include <cassert>

int main() {
    int x = 10;
    assert(x > 0);
}
\end{lstlisting}
条件为假时，程序会在运行时终止。因此，\texttt{assert} 属于运行期检查。

\section{什么是 \texttt{static\_assert}}
\texttt{static\_assert} 用于在编译阶段检查条件。
\begin{lstlisting}
static_assert(条件, "错误信息");
\end{lstlisting}
例如：
\begin{lstlisting}
static_assert(sizeof(int) >= 4,
              "int 类型至少需要 4 字节");
\end{lstlisting}
条件为假时，程序无法通过编译。

\begin{conceptbox}{核心区别}
\begin{itemize}
    \item \texttt{assert}：运行时检查；
    \item \texttt{static\_assert}：编译时检查。
\end{itemize}
\end{conceptbox}

\section{条件必须能在编译期确定}
正确：
\begin{lstlisting}
constexpr int size = 10;
static_assert(size > 0, "size 必须大于 0");
\end{lstlisting}
错误：
\begin{lstlisting}
int size;
std::cin >> size;
static_assert(size > 0, "size 必须大于 0");
\end{lstlisting}
第二段中的 \texttt{size} 只有运行后才能确定，因此不能用于 \texttt{static\_assert}。

\section{C++17 的简化写法}
在较早标准中通常需要写错误信息：
\begin{lstlisting}
static_assert(sizeof(int) >= 4, "int 太小");
\end{lstlisting}
从 C++17 开始，可以省略第二个参数：
\begin{lstlisting}
static_assert(sizeof(int) >= 4);
\end{lstlisting}
不过实际开发中，保留错误信息更方便定位问题。

\section{常见用途}
\subsection{检查类型大小}
\begin{lstlisting}
static_assert(sizeof(long long) >= 8,
              "long long 至少需要 8 字节");
\end{lstlisting}

\subsection{检查模板参数}
\begin{lstlisting}
template <typename T>
T add(T a, T b) {
    static_assert(sizeof(T) >= 4,
                  "类型长度不能小于 4 字节");
    return a + b;
}
\end{lstlisting}

\subsection{结合类型特征}
\begin{lstlisting}
#include <type_traits>

template <typename T>
T twice(T value) {
    static_assert(std::is_integral_v<T>,
                  "T 必须是整数类型");
    return value * 2;
}
\end{lstlisting}
调用：
\begin{lstlisting}
twice(10);     // 正确
twice(3.14);   // 编译错误
\end{lstlisting}

\subsection{检查数组长度}
\begin{lstlisting}
template <typename T, std::size_t N>
void process(T (&arr)[N]) {
    static_assert(N >= 5,
                  "数组长度至少为 5");
}
\end{lstlisting}

\subsection{检查配置常量}
\begin{lstlisting}
constexpr int BufferSize = 1024;

static_assert(BufferSize > 0,
              "缓冲区大小必须大于 0");

static_assert((BufferSize & (BufferSize - 1)) == 0,
              "缓冲区大小必须是 2 的幂");
\end{lstlisting}


\section{与 \texttt{assert} 对比}
\begin{center}
\begin{tabular}{lll}
\toprule
项目 & \texttt{assert} & \texttt{static\_assert} \\
\midrule
检查时间 & 运行时 & 编译时 \\
条件要求 & 普通布尔表达式 & 常量表达式 \\
失败结果 & 程序终止 & 编译失败 \\
是否需要头文件 & \texttt{<cassert>} & 不需要 \\
典型用途 & 检查运行数据 & 检查类型和配置 \\
\bottomrule
\end{tabular}
\end{center}

\section{常见错误}
\subsection{使用运行时变量}
\begin{lstlisting}
int x = 10;
static_assert(x > 0, "x 必须大于 0"); // 错误
\end{lstlisting}
如果要检查运行时数据，应使用普通判断或 \texttt{assert}。

\subsection{条件和错误信息矛盾}
错误：
\begin{lstlisting}
static_assert(sizeof(int) < 4,
              "int 至少应有 4 字节");
\end{lstlisting}
正确：
\begin{lstlisting}
static_assert(sizeof(int) >= 4,
              "int 至少应有 4 字节");
\end{lstlisting}



\section{总结}
\begin{conceptbox}{板书式记忆}
\begin{lstlisting}
static_assert(编译期条件, "错误信息");
\end{lstlisting}
\begin{enumerate}[label=\arabic*.]
    \item 条件必须在编译期确定；
    \item 条件为假时，程序无法通过编译；
    \item 常用于检查类型、模板参数和编译期配置。
\end{enumerate}
\end{conceptbox}

\end{document}
